\documentclass[12pt,letterpaper]{article}

\usepackage[utf8]{inputenc}
\usepackage[T1]{fontenc}
\usepackage[spanish,es-nodecimaldot,es-noshorthands]{babel}
\usepackage{geometry}
\usepackage{amsmath}
\usepackage{tikz}
\usepackage{xcolor}
\usepackage{float}
\usepackage{enumitem}
\usepackage{titlesec}

\usetikzlibrary{arrows.meta, positioning, shapes.geometric}

\geometry{left=2.3cm,right=2.3cm,top=2.2cm,bottom=2.2cm}

\titleformat{\section}{\large\bfseries}{}{0em}{}
\setlength{\parindent}{0pt}
\setlength{\parskip}{0.5em}

\begin{document}

\begin{center}
    {\LARGE \textbf{Proyecto 1: AGV para logística de alta carga}}\\[0.3cm]
    {\large Propuesta corta de solución}
\end{center}

\section{Propuesta}

Para este AGV, la idea es que el vehículo transporte hasta \textbf{300 kg}, por lo que necesita \textbf{motores de alto torque} y desplazarse a \textbf{baja velocidad} para moverse con seguridad en pasillos estrechos.

La propuesta es dividir el sistema en dos partes:

\begin{itemize}[leftmargin=1.2cm]
    \item \textbf{ESP32:} encargado de controlar los motores, es decir, decidir cuándo avanzar, detenerse o reducir velocidad.
    \item \textbf{Raspberry Pi:} encargada de procesar la información de una \textbf{cámara} y de un \textbf{sensor infrarrojo}, para detectar personas u obstáculos.
\end{itemize}

La Raspberry Pi analiza en todo momento si hay una persona cerca. Si detecta riesgo, envía por \textbf{MQTT} una señal al ESP32 para que el AGV se detenga o reduzca velocidad.

Además, el sensor infrarrojo puede actuar como respaldo a una distancia aproximada de \textbf{1 a 2 metros}, para evitar que el vehículo frene demasiado tarde. Esto es importante porque, si frena de golpe, la carga puede desplazarse hacia adelante por \textbf{inercia}.

\section{Hardware y comunicación}

\begin{itemize}[leftmargin=1.2cm]
    \item \textbf{Motores DC con alto torque o motorreductores:} para mover la carga de 300 kg.
    \item \textbf{Driver de motores:} para alimentar y controlar los motores.
    \item \textbf{ESP32:} control directo de movimiento.
    \item \textbf{Raspberry Pi:} procesamiento de visión y lógica de detección.
    \item \textbf{Cámara:} detección visual de personas.
    \item \textbf{Sensor infrarrojo:} detección rápida a corta distancia.
    \item \textbf{MQTT:} protocolo para comunicar Raspberry Pi y ESP32.
\end{itemize}

\section{Diagrama general del sistema}

\begin{figure}[H]
\centering
\begin{tikzpicture}[
    node distance=1cm and 1.4cm,
    bloque/.style={
        rectangle,
        rounded corners,
        draw,
        thick,
        minimum width=3.6cm,
        minimum height=1cm,
        align=center,
        fill=blue!10
    },
    sensor/.style={
        rectangle,
        rounded corners,
        draw,
        thick,
        minimum width=3.6cm,
        minimum height=1cm,
        align=center,
        fill=green!10
    },
    control/.style={
        rectangle,
        rounded corners,
        draw,
        thick,
        minimum width=3.8cm,
        minimum height=1cm,
        align=center,
        fill=orange!15
    },
    flecha/.style={
        -{Latex[length=2.5mm]},
        thick
    }
]

\node[sensor] (cam) {Cámara};
\node[sensor, below=of cam] (ir) {Sensor infrarrojo\\(1--2 m)};
\node[bloque, right=2cm of ir] (rasp) {Raspberry Pi\\Procesa detección};
\node[control, right=2cm of rasp] (esp) {ESP32\\Control de movimiento};
\node[bloque, below=of esp] (driver) {Driver de motores};
\node[bloque, below=of driver] (motor) {Motores de alto torque};
\node[bloque, below=of motor] (agv) {AGV con carga de 300 kg};

\draw[flecha] (cam) -- (rasp);
\draw[flecha] (ir) -- (rasp);
\draw[flecha] (rasp) -- node[above]{MQTT} (esp);
\draw[flecha] (esp) -- (driver);
\draw[flecha] (driver) -- (motor);
\draw[flecha] (motor) -- (agv);

\end{tikzpicture}
\caption{Arquitectura básica propuesta para el AGV.}
\end{figure}

\section{Diagrama de funcionamiento}

\begin{figure}[H]
\centering
\begin{tikzpicture}[
    node distance=0.8cm,
    iniciofin/.style={
        ellipse,
        draw,
        thick,
        minimum width=3.5cm,
        minimum height=1cm,
        align=center,
        fill=green!15
    },
    proceso/.style={
        rectangle,
        rounded corners,
        draw,
        thick,
        minimum width=5.2cm,
        minimum height=1cm,
        align=center,
        fill=blue!10
    },
    decision/.style={
        diamond,
        draw,
        thick,
        aspect=2.2,
        minimum width=3.8cm,
        align=center,
        fill=orange!15
    },
    flecha/.style={
        -{Latex[length=2.5mm]},
        thick
    }
]

\node[iniciofin] (inicio) {Inicio};
\node[proceso, below=of inicio] (avance) {ESP32 mueve el AGV\\a baja velocidad};
\node[proceso, below=of avance] (monitorea) {Cámara y sensor IR\\monitorean el entorno};
\node[decision, below=of monitorea] (persona) {¿Se detecta\\persona u obstáculo?};
\node[proceso, right=3.3cm of persona] (senal) {Raspberry Pi envía\\señal MQTT al ESP32};
\node[proceso, below=of persona] (seguir) {El AGV sigue avanzando};
\node[proceso, below=of senal] (paro) {ESP32 detiene o reduce\\la velocidad de los motores};
\node[iniciofin, below=of seguir] (fin) {Fin / ciclo continuo};

\draw[flecha] (inicio) -- (avance);
\draw[flecha] (avance) -- (monitorea);
\draw[flecha] (monitorea) -- (persona);
\draw[flecha] (persona) -- node[above]{Sí} (senal);
\draw[flecha] (senal) -- (paro);
\draw[flecha] (persona) -- node[right]{No} (seguir);
\draw[flecha] (seguir) -- (fin);
\draw[flecha] (paro) |- (fin);

\end{tikzpicture}
\caption{Flujo básico de operación del AGV.}
\end{figure}

\section{Conclusión}

En resumen, esta propuesta usa un \textbf{ESP32} para el control de motores y una \textbf{Raspberry Pi} para procesar la detección de personas mediante cámara y sensor infrarrojo. La comunicación entre ambos se realiza con \textbf{MQTT}. Como el AGV transporta \textbf{300 kg}, debe usar motores de alto torque y avanzar lentamente para evitar riesgos, sobre todo en pasillos estrechos y cuando haya personas cerca.

\end{document}
```
